\subsubsection*{(ii) Trouver une méthode de construction qui se rapproche de $m(\ell)$.}

Pour approcher la valeur de $m(\ell)$, nous avons développé une méthode de construction. Elle consiste à utiliser pour un $1 \leq n \leq \ell$ fixé:

\begin{itemize}
	\item Les $n$ premières graduations: $0, 1, 2, \ldots, n-1$ $\implies n$ graduations;
	\item Tous les multiples suivants de $n$ soustrait de 1: $2n-1, 3n-1, \ldots$, jusqu'à la plus grande valeur inférieure à $\ell$ $\implies \left\lfloor \frac{\ell}{n} \right\rfloor - 1$ graduations;
	\item Le nombre $\ell$ lui-même $\implies$ 1 graduation.
\end{itemize}

Cette solution est valide car:
\begin{itemize}
	\item Avec les $n$ premières graduations, on peut mesurer toutes les distances de 1 à $n-1$.
	\item Pour chaque multiple de $n$ soustrait de 1, on peut mesurer toutes les distances de $(k-1)n$ à $kn - 1$ à l'aide des $n$ premières graduations.
	\item Avec la graduation à $\ell$, on peut mesurer toutes les distances restantes jusqu'à $\ell$.
\end{itemize}

Ainsi, le nombre total de graduations utilisées pour un $n$ fixé est:
\[
n + \left\lfloor \frac{\ell}{n} \right\rfloor - 1 + 1 = n + \left\lfloor \frac{\ell}{n} \right\rfloor.
\]

En dérivant cette valeur par rapport à $n$, nous remarquons que le nombre de graduations est minimisé lorsque $n$ est proche de $\sqrt{\ell}$. En effet, en posant $n = \lceil \sqrt{\ell} \rceil$, nous obtenons un nombre de graduations d'environ:

\[
\lceil \sqrt{\ell} \rceil + \left\lfloor \frac{\ell}{\lceil \sqrt{\ell} \rceil} \right\rfloor \approx 2\sqrt{\ell}.
\]

Cette méthode de construction nous permet donc d'affirmer que $m(\ell) \lesssim 2\sqrt{\ell}$.